\documentclass[../../main.tex]{subfiles}
\begin{document}
\section{横截性}\label{section:ch06-横截性}

作为 Sard 定理的最后一个应用,我们说明如何扰动子流形使它们``良好地''相交。为解释这意味着什么,我们引进横截性的概念。

向量空间的两个线性子空间之交总是另一线性子空间。对子流形而言,相应的命题当然不成立:容易举出光滑子流形之交不是子流形的例子(见习题 6-14)。但对子流形附加一条假设后,即可证明它们的交仍是子流形。

设 $M$ 是光滑流形。称 $M$ 的两个嵌入子流形 $S,S'\subseteq M$ 横截相交,如果对每一点 $p\in S\cap S'$,切空间 $T_pS$ 与 $T_pS'$ 张成 $T_pM$(其中我们把 $T_pS$ 与 $T_pS'$ 视为 $T_pM$ 的子空间)。

对许多目的,使用下面更一般的定义更为方便。若 $F:N\to M$ 是光滑映射, $S\subseteq M$ 是嵌入子流形,我们说 $F$ 横截于 $S$,如果对每一点 $x\in F^{-1}(S)$,空间 $T_{F(x)}S$ 与 $\mathrm{d}F_x(T_xN)$ 张成 $T_{F(x)}M$。一个特例值得注意:若 $F$ 是光滑淹没,则它自动横截于 $M$ 的每个嵌入子流形。两个嵌入子流形横截相交,当且仅当其中任一个到另一个的包含映射横截于另一个。

下一个结果是正则水平集定理的推广,它说明横截性为何可取。

\begin{theorem}\label{theorem:ch06-transverse}
设 $N$ 与 $M$ 是光滑流形, $S\subseteq M$ 是嵌入子流形。
\begin{enumerate}[(a)]
\item 若 $F:N\to M$ 是横截于 $S$ 的光滑映射,则 $F^{-1}(S)$ 是 $N$ 的嵌入子流形,其余维数等于 $S$ 在 $M$ 中的余维数。
\item 若 $S'\subseteq M$ 是与 $S$ 横截相交的嵌入子流形,则 $S\cap S'$ 是 $M$ 的嵌入子流形,其余维数等于 $S$ 与 $S'$ 的余维数之和。
\end{enumerate}
\end{theorem}

\begin{proof}
第二个命题容易由第一个推出,只需取 $F$ 为包含映射 $S'\hookrightarrow M$;并注意光滑嵌入之复合 $S\cap S'\hookrightarrow S\hookrightarrow M$ 仍是光滑嵌入。

为证 (a),令 $m$ 表示 $M$ 的维数, $k$ 表示 $S$ 在 $M$ 中的余维数。给定 $x\in F^{-1}(S)$,我们可在 $M$ 中找到 $F(x)$ 的邻域 $U$ 与 $S$ 的局部定义函数 $\varphi:U\to\mathbb{R}^k$,使 $S\cap U=\varphi^{-1}(0)$。若能证明 $0$ 是 $\varphi\circ F$ 的正则值,定理即证,因为 $F^{-1}(S)\cap F^{-1}(U)$ 是 $\varphi\circ F|_{F^{-1}(U)}$ 的零点集。

给定 $z\in T_0\mathbb{R}^k$ 与 $p\in(\varphi\circ F)^{-1}(0)$, $0$ 是 $\varphi$ 的正则值这一事实意味着存在向量 $y\in T_{F(p)}M$ 使 $\mathrm{d}\varphi_{F(p)}(y)=z$。$F$ 横截于 $S$ 这一事实意味着可写 $y=y_0+\mathrm{d}F_p(v)$,其中 $y_0\in T_{F(p)}S$, $v\in T_pN$。由于 $\varphi$ 在 $S\cap U$ 上为常数,故 $\mathrm{d}\varphi_{F(p)}(y_0)=0$,于是
\begin{align*}
\mathrm{d}(\varphi\circ F)_p(v)=\mathrm{d}\varphi_{F(p)}\bigl(\mathrm{d}F_p(v)\bigr)
=\mathrm{d}\varphi_{F(p)}\bigl(y_0+\mathrm{d}F_p(v)\bigr)
=\mathrm{d}\varphi_{F(p)}(y)=z.
\end{align*}
于是 $F^{-1}(S)$ 是余维数 $k$ 的嵌入子流形。
\end{proof}

例如,在 $\mathbb{R}^3$ 中,该定理表明横截相交的光滑曲线与光滑曲面只在孤立点相交,而两个横截相交的光滑曲面横截相交于一条光滑曲线。$\mathbb{R}^3$ 中两条光滑曲线横截相交,当且仅当其交为空,因为在任一交点处,两条曲线的两个一维切空间无法张成 $T\mathbb{R}^3$。

由于淹没横截于每个嵌入子流形,下面的推论是直接的。

\begin{corollary}\label{corollary:ch06-submersion-preimage-transverse}
设 $N$ 与 $M$ 是光滑流形, $S\subseteq M$ 是余维数 $k$ 的嵌入子流形, $F:N\to M$ 是淹没。则 $F^{-1}(S)$ 是 $N$ 的余维数 $k$ 嵌入子流形。
\end{corollary}

横截性还提供了把一个子流形识别为图的方便判别法。下一定理是隐函数定理的整体版本。

\begin{theorem}[Global Characterization of Graphs]\label{theorem:ch06-global-graph}
设 $M$ 与 $N$ 是光滑流形, $S\subseteq M\times N$ 是浸入子流形。令 $\pi_M$ 与 $\pi_N$ 分别表示 $M\times N$ 到 $M$ 与到 $N$ 的投影。下列条件等价。
\begin{enumerate}[(a)]
\item $S$ 是某光滑映射 $f:M\to N$ 的图。
\item $\pi_M|_S$ 是从 $S$ 到 $M$ 上的微分同胚。
\item 对每个 $p\in M$,子流形 $S$ 与 $\{p\}\times N$ 在恰一点处横截相交。
\end{enumerate}
若这些条件成立,则 $S$ 是由 $f=\pi_N\circ(\pi_M|_S)^{-1}$ 定义的映射 $f:M\to N$ 的图。
\end{theorem}

\begin{proof}
见习题 6-15。
\end{proof}

\begin{corollary}[Local Characterization of Graphs]\label{corollary:ch06-local-graph}
设 $M$ 与 $N$ 是光滑流形, $S\subseteq M\times N$ 是浸入子流形,且 $(p,q)\in S$。若 $S$ 在 $(p,q)$ 处横截于子流形 $\{p\}\times N$,则存在 $p$ 在 $M$ 中的邻域 $U$ 与 $(p,q)$ 在 $S$ 中的邻域 $V$,使 $V$ 是某光滑映射 $f:U\to N$ 的图。
\end{corollary}

\begin{proof}
假设保证 $\mathrm{d}(\pi_M)_{(p,q)}:T_{(p,q)}S\to T_pM$ 是同构,故 $\pi_M|_S$ 限制为从 $S$ 中 $(p,q)$ 的某邻域 $V$ 到 $p$ 的某邻域 $U$ 上的微分同胚。结论由\cref{theorem:ch06-global-graph}(b) 推出。
\end{proof}

横截相交子流形与横截映射令人惊讶之处在于它们是``一般的'',正如我们即将看到的。为铺垫,我们需要考虑比光滑同伦稍一般的映射族。设 $N,M,S$ 是光滑流形,并对每个 $s\in S$ 给定映射 $F_s:N\to M$。若由 $F(x,s)=F_s(x)$ 定义的映射 $F:N\times S\to M$ 光滑,则称族 $\{F_s:s\in S\}$ 为\textbf{光滑映射族}。你可把这样的族视为同伦的高维类比。下一个命题说明这种族与普通同伦的关系。

\begin{proposition}\label{proposition:ch06-smooth-family-homotopic}
若 $\{F_s:s\in S\}$ 是从 $N$ 到 $M$ 的光滑映射族,且 $S$ 连通,则对任意 $s_1,s_2\in S$,映射 $F_{s_1},F_{s_2}:N\to M$ 同伦。
\end{proposition}

\begin{proof}
由于 $S$ 连通,它道路连通。若 $\gamma:[0,1]\to S$ 是从 $s_1$ 到 $s_2$ 的任一路径,则 $H(x,s)=F\bigl(x,\gamma(s)\bigr)$ 是从 $F_{s_1}$ 到 $F_{s_2}$ 的同伦。
\end{proof}

寻找横截映射的关键是下面 Sard 定理的应用,它给出一族光滑映射中至少含一个横截于给定子流形的映射这一简单充分条件。若 $S$ 是光滑流形, $B\subseteq S$ 是补集在 $S$ 中零测的子集,我们说 $B$ 包含 $S$ 的几乎每个元素。

\begin{theorem}[Parametric Transversality Theorem]\label{theorem:ch06-parametric-transversality}
设 $N$ 与 $M$ 是光滑流形, $X\subseteq M$ 是嵌入子流形, $\{F_s:s\in S\}$ 是从 $N$ 到 $M$ 的光滑映射族。若映射 $F:N\times S\to M$ 横截于 $X$,则对几乎每个 $s\in S$,映射 $F_s:N\to M$ 横截于 $X$。
\end{theorem}

\begin{proof}
假设蕴涵 $W=F^{-1}(X)$ 是 $N\times S$ 的嵌入子流形(由\cref{theorem:ch06-transverse})。令 $\pi:N\times S\to S$ 为到第二分量的投影。我们实际要证明的是:若 $s\in S$ 是限制 $\pi|_W$ 的正则值,则 $F_s$ 横截于 $X$。由于由 Sard 定理几乎每个 $s$ 都是正则值,这就证明了定理。

设 $s\in S$ 是 $\pi|_W$ 的正则值。任取 $p\in F_s^{-1}(X)$,并令 $q=F_s(p)\in X$。我们需证明 $T_qM=T_qX+\mathrm{d}(F_s)(T_pN)$。已知如下。首先,由于我们对 $F$ 的假设,
\begin{align}\label{align:ch06-tm-decomp1}
T_qM=T_qX+\mathrm{d}F\bigl(T_{(p,s)}(N\times S)\bigr).
\end{align}
其次,由于 $s$ 是正则值且 $(p,s)\in W$,
\begin{align}\label{align:ch06-ts-decomp}
T_sS=\mathrm{d}\pi\bigl(T_{(p,s)}W\bigr).
\end{align}
第三,由习题 6-10 的结果, $T_{(p,s)}W=\bigl(\mathrm{d}F_{(p,s)}\bigr)^{-1}(T_qX)$,这蕴涵
\begin{align}\label{align:ch06-dfw}
\mathrm{d}F\bigl(T_{(p,s)}W\bigr)=T_qX.
\end{align}

现令 $w\in T_qM$ 任意。我们需找 $v\in T_qX$ 与 $y\in T_pN$,使
\begin{align}\label{align:ch06-w-decomp}
w=v+\mathrm{d}(F_s)(y).
\end{align}
由\eqref{align:ch06-tm-decomp1},存在 $v_1\in T_qX$ 与 $(y_1,z_1)\in T_pN\times T_sS\cong T_{(p,s)}(N\times S)$,使
\begin{align}\label{align:ch06-w-decomp2}
w=v_1+\mathrm{d}F(y_1,z_1).
\end{align}
由\eqref{align:ch06-ts-decomp},存在 $(y_2,z_2)\in T_{(p,s)}W$ 使 $\mathrm{d}\pi(y_2,z_2)=z_1$。由于 $\pi$ 是投影,这意味着 $z_2=z_1$。由线性性,可写
\begin{align*}
\mathrm{d}F(y_1,z_1)=\mathrm{d}F(y_2,z_1)+\mathrm{d}F(y_1-y_2,0).
\end{align*}
一方面, \eqref{align:ch06-dfw} 蕴涵
\begin{align*}
\mathrm{d}F(y_2,z_1)=\mathrm{d}F(y_2,z_2)\in\mathrm{d}F\bigl(T_{(p,s)}W\bigr)=T_qX.
\end{align*}
另一方面,若 $\sigma_s:N\to N\times S$ 是映射 $\sigma_s(p_0)=(p_0,s)$,则有 $F_s=F\circ\sigma_s$ 且 $\mathrm{d}(\sigma_s)(y_1-y_2)=(y_1-y_2,0)$,因此
\begin{align*}
\mathrm{d}F(y_1-y_2,0)=\mathrm{d}F\circ\mathrm{d}(\sigma_s)(y_1-y_2)=\mathrm{d}(F_s)(y_1-y_2).
\end{align*}
于是由\eqref{align:ch06-w-decomp2},取 $v=v_1+\mathrm{d}F(y_2,z_1)$ 与 $y=y_1-y_2$ 即满足\eqref{align:ch06-w-decomp},证毕。
\end{proof}

为应用参数横截性定理,我们需要构造满足假设的光滑映射族。下一定理的证明表明总能做到。

\begin{theorem}[Transversality Homotopy Theorem]\label{theorem:ch06-transversality-homotopy}
设 $M$ 与 $N$ 是光滑流形, $X\subseteq M$ 是嵌入子流形。每个光滑映射 $f:N\to M$ 都同伦于一个横截于 $X$ 的光滑映射 $g:N\to M$。
\end{theorem}

\begin{proof}
证明的关键是构造光滑映射 $F:N\times S\to M$ 横截于 $X$,其中 $S=B^k$ 对某 $k$,且 $F_0=f$。于是由参数横截性定理,存在某 $s\in S$ 使 $F_s:N\to M$ 横截于 $X$;并由\cref{proposition:ch06-smooth-family-homotopic} 知 $F_s$ 同伦于 $f$。

由 Whitney 嵌入定理,我们可设 $M$ 是某 $\mathbb{R}^k$ 的真嵌入子流形。令 $U$ 为 $M$ 在 $\mathbb{R}^k$ 中的管状邻域,并令 $r:U\to M$ 为光滑收缩映射,它同时是光滑淹没。若由\eqref{align:ch06-delta-tube} 定义 $\delta:M\to\mathbb{R}^+$,则\cref{corollary:ch06-small-positive-function} 表明存在光滑函数 $\varepsilon:N\to\mathbb{R}^+$,处处满足 $0<\varepsilon(p)<\delta(f(p))$。令 $S$ 为 $\mathbb{R}^k$ 中的单位球,并定义 $F:N\times S\to M$ 为
\begin{align*}
F(p,s)=r\bigl(f(p)+\varepsilon(p)s\bigr).
\end{align*}
注意 $|\varepsilon(p)s|<\varepsilon(p)<\delta(f(p))$,这蕴涵 $f(p)+\varepsilon(p)s\in U$,故 $F$ 良定义。显然 $F$ 光滑,且 $F_0=f$,因为 $r$ 是收缩映射。

对每个 $p\in N$, $F$ 在 $\{p\}\times S$ 上的限制是局部微分同胚 $s\mapsto f(p)+\varepsilon(p)s$ 接着光滑淹没 $r$ 的复合,故 $F$ 是光滑淹没,从而横截于 $X$。
\end{proof}

\end{document}